\documentclass[12pt,a4paper]{article}

\usepackage[a4paper,margin=1in]{geometry}
\usepackage[T1]{fontenc}
\usepackage[utf8]{inputenc}
\usepackage{enumitem}
\setlength{\emergencystretch}{3em}

\title{\textbf{Supervisor Meeting Minutes}\\Meeting 7}
\author{Present: Matthew Chalmers, Jonas Leipner}
\date{Date: 17 November 2026}

\begin{document}
\maketitle

\section{Achieved in Previous Weeks}

\noindent\textit{Previous weeks: 3--17 November 2026 (no meeting was held on
10 November)}

\begin{itemize}[leftmargin=*]
    \item Rebuilt the game from scratch on the main branch; the prototype is
    kept for archival purposes only.
    \item Separated the game into a model, an engine and a view, so that Flame
    only draws the game and passes on player input.
    \item Implemented tower placement with the limit of five towers (FR-M1)
    and the score for defeated enemies (FR-S1).
    \item Moved the combat logic into a \texttt{dart\_eval} script, with the
    native implementation as a fallback when a script fails (FR-M2).
    \item Made it possible to swap the combat script while the game is
    running, without restarting the application.
    \item Recorded the architecture decisions in the decision log and set up
    automated tests that run on every push.
\end{itemize}

\section{Issues Discussed \& Solutions Agreed}

\begin{itemize}[leftmargin=*]
    \item \textbf{Demonstration:} The working game was demonstrated, including
    swapping the combat script during play so that towers change from
    targeting the leading enemy to the strongest enemy.
    \begin{itemize}
        \item \textbf{Solution:} The current base is suitable to build the
        research on.
    \end{itemize}
    \item \textbf{Limitations:} The current limitations were discussed:
    \texttt{dart\_eval} has no time limit, so a script that never returns
    freezes the game; compiling a script can briefly stall the game; and the
    scripts are still bundled with the application rather than downloaded.
    \begin{itemize}
        \item \textbf{Solution:} These are recorded and do not block the next
        steps.
    \end{itemize}
    \item \textbf{Project direction:} Four directions were compared: the
    technical cost of the sandbox (A), the safety of dynamically loaded logic
    (B), shared content between players (C), and live balancing (D).
    \begin{itemize}
        \item \textbf{Solution:} A and B are more standard and less
        interesting as research, and can be done at any point. The project
        will concentrate on C and D.
    \end{itemize}
    \item \textbf{Participants:} It is not yet known who is available to
    interview and to play test the game, and whether they would be developers
    willing to write custom code or players who would only play.
    \begin{itemize}
        \item \textbf{Solution:} Find out who is available and which group
        they belong to, as this decides which features are worth building.
    \end{itemize}
    \item \textbf{Incubator for non-developers:} The incubator was discussed as
    a way for players who do not write code to still create new towers.
    \begin{itemize}
        \item \textbf{Solution:} Research games that already use genetic or
        evolving mechanics, such as Skylanders, the Game of Life, and games
        where species or worlds mutate and evolve, and use them as examples to
        build on.
    \end{itemize}
    \item \textbf{Interviews and ethics:} Early interviews are planned to
    guide the game and the research. After a short, non-technical explanation
    of dynamic loading and of the incubator, participants will be asked which
    aesthetic they prefer (cute or more intense and edgy), whether the
    incubator sounds appealing, and what they would find fun.
    \begin{itemize}
        \item \textbf{Solution:} Obtain the ethics checklist and have
        participants sign it before the interviews.
    \end{itemize}
\end{itemize}

\section{Objectives for This Week}

\noindent\textit{Upcoming week: 17--24 November 2026}

\begin{itemize}[leftmargin=*]
    \item Find out who is available to interview and play test, and whether
    they are developers or players.
    \item Research games with genetic, mutating or evolving mechanics as
    examples for the incubator.
    \item Obtain the ethics checklist.
    \item Prepare the interview questions and a simple explanation of dynamic
    loading and the incubator for non-developers.
\end{itemize}

\section*{\small AI-use Disclaimer}
\small The meeting content and source notes for these minutes are genuine.
Artificial intelligence was used only to organise and summarise the notes
and transcript into a more digestible format; it was not used to invent the
meeting content or decisions.

\end{document}
